% Part: first-order-logic
% Chapter: syntax-and-semantics
% Section: unique-readability

\documentclass[../../../include/open-logic-section]{subfiles}

\begin{document}

\olfileid{fol}{syn}{unq}

\olsection{خوانش یکتا}

\begin{explain}
شیوه‌ای که !!{formula}s را تعریف کردیم تضمین می‌کند که هر !!{formula}
\emph{خوانشی یکتا} دارد؛ یعنی در اصل فقط یک راه برای ساختن آن طبق قواعد
ساخت !!{formula}s و تنها یک راه برای «تفسیر» آن وجود دارد. اگر چنین نبود،
!!{formula}sی مبهم می‌داشتیم، یعنی !!{formula}sی با بیش از یک خوانش یا
تفسیر---و آشکار است که می‌خواهیم از این امر پرهیز کنیم. مهم‌تر آنکه بدون
این ویژگی، بیشتر تعریف‌ها و برهان‌هایی که ارائه خواهیم کرد به نتیجه
نمی‌رسیدند.

شاید بهترین راه روشن‌کردن این نکته آن باشد که ببینیم اگر قواعد نادرستی برای
ساخت !!{formula}s می‌دادیم که خوانش یکتا را تضمین نمی‌کردند، چه پیش
می‌آمد. برای نمونه، ممکن بود پرانتزها را در قواعد ساخت رابط‌ها فراموش کنیم
و مثلاً قاعدهٔ زیر را مجاز بدانیم:
\begin{quote}
اگر $!A$ و $!B$ از !!{formula}s باشند، آنگاه $!A \lif !B$ نیز چنین است.
\end{quote}
اگر از فرمول اتمی $!D$ آغاز کنیم، این قاعده اجازه می‌دهد $!D \lif !D$ را
بسازیم. سپس از این فرمول همراه با $!D$، فرمول $!D \lif !D \lif !D$ را
به دست می‌آوریم. اما این کار را به دو شیوه می‌توان انجام داد:
\begin{enumerate}
\item $!D$ را $!A$ و $!D \lif !D$ را $!B$ می‌گیریم.
\item $!A$ را $!D \lif !D$ و $!B$ را~$!D$ می‌گیریم.
\end{enumerate}
به همین ترتیب، برای «خواندن» !!{formula}~$!D \lif !D \lif !D$ نیز دو
راه وجود دارد. این فرمول به صورت $!B \lif !C$ است که در آن $!B$ همان
$!D$ و $!C$ همان $!D \lif !D$ است؛ اما \emph{همچنین} به صورت
$!B \lif !C$ است که در آن $!B$ برابر $!D \lif !D$ و $!C$ برابر~$!D$ است.

اگر چنین اتفاقی بیفتد، تعریف‌های ما همیشه درست کار نخواهند کرد. برای نمونه،
هنگام تعریف !!{main operator} یک فرمول می‌گوییم: در فرمولی به صورت
$!B \lif !C$، !!{main operator} همان رخداد نشان‌داده‌شدهٔ~$\lif$ است. اما
اگر بتوانیم فرمول $!D \lif !D \lif !D$ را به دو شیوهٔ یادشده با
$!B \lif !C$ تطبیق دهیم، در یک حالت رخداد نخست $\lif$ و در حالت دیگر
رخداد دوم آن را به‌عنوان !!{main operator} به دست می‌آوریم. حال آنکه
می‌خواهیم !!{main operator} \emph{تابعی} از !!{formula} باشد؛ یعنی هر
!!{formula} باید دقیقاً یک رخداد !!{main operator} داشته باشد.
\end{explain}

\begin{lem}
شمار پرانتزهای باز و بسته در !!a{formula}~$!A$ برابر است.
\end{lem}

\begin{proof}
این ادعا را با استقرا بر شیوهٔ ساخت $!A$ اثبات می‌کنیم. دو کار لازم است:
(الف) نخست باید ثابت کنیم همهٔ فرمول‌های اتمی ویژگی مورد نظر را دارند
(پایهٔ استقرا). (ب) سپس باید ثابت کنیم هرگاه از فرمول‌های داده‌شده
فرمول‌هایی تازه بسازیم، اگر فرمول‌های پیشین ویژگی را داشته باشند،
فرمول‌های تازه نیز آن را دارند.

بگذارید $l(!A)$ شمار پرانتزهای باز و $r(!A)$ شمار پرانتزهای بسته در~$!A$
باشد؛ به همین ترتیب، $l(t)$ و $r(t)$ به‌ترتیب شمار پرانتزهای باز و بسته در
ترم~$t$ باشند.

\begin{prob}
    اثبات کنید برای هر ترم~$t$، $l(t) = r(t)$.
\end{prob}

\begin{enumerate}
\tagitem{prvFalse}{\indcase{!A}{\lfalse}{$\indfrm$ دارای $0$ پرانتز باز و
    $0$ پرانتز بسته است.}}{}

\tagitem{prvTrue}{\indcase{!A}{\ltrue}{$\indfrm$ دارای $0$ پرانتز باز و
    $0$ پرانتز بسته است.}}{}

\item \indcase{!A}{\Atom{R}{t_1,\dots,t_n}}{$l(\indfrm) = 1 + l(t_1) +
  \dots + l(t_n) = 1 + r(t_1) + \dots + r(t_n) = r(\indfrm)$. در اینجا
  از این واقعیت استفاده می‌کنیم که، چنان‌که در تمرین خواسته شد، برای هر
  ترم~$t$ داریم $l(t) = r(t)$.}

\item \indcase{!A}{\eq[t_1][t_2]}{$l(\indfrm) = l(t_1) + l(t_2) =
  r(t_1) + r(t_2) = r(\indfrm)$.}

\tagitem{prvNot}{\indcase{!A}{\lnot !B}{بنا بر فرض استقرا،
    $l(!B) = r(!B)$. پس $l(\indfrm) = l(!B) = r(!B) =
    r(\indfrm)$.}}{}

\item \indcase{!A}{(!B \ast !C)}{بنا بر فرض استقرا، $l(!B) =
  r(!B)$ و $l(!C) = r(!C)$. پس $l(\indfrm) = 1 + l(!B) + l(!C) =
  1 + r(!B) + r(!C) = r(\indfrm)$.}

\tagitem{prvAll}{\indcase{!A}{\lforall[x][!B]}{بنا بر فرض استقرا،
    $l(!B) = r(!B)$. پس $l(\indfrm) = l(!B) = r(!B) =
    r(\indfrm)$.}}{}

% Print case for \lexists if prvEx and not prvAll
\tagitem{prvAll,notprvEx}{}{\indcase{!A}{\lexists[x][!B]}{بنا بر فرض
    استقرا، $l(!B) = r(!B)$. پس $l(\indfrm) = l(!B) = r(!B) =
    r(\indfrm)$.}}

% Just say similarly if prvEx and prvAll
\tagitem{notprvAll,notprvEx}{}{\indcase{!A}{\lexists[x][!B]}{به همین ترتیب.}}
\end{enumerate}
\end{proof}

\begin{defn}[پیشوند سره]
رشتهٔ نمادهای $!B$ \emph{پیشوند سرهٔ} رشتهٔ نمادهای~$!A$ است هرگاه الحاق
$!B$ با رشته‌ای ناتهی از نمادها،~$!A$ را به دست دهد.
\end{defn}

\begin{lem}\ollabel{lem:no-prefix}
اگر $!A$ یک !!a{formula} و $!B$ پیشوند سرهٔ $!A$ باشد، آنگاه $!B$ یک
!!a{formula} نیست.
\end{lem}

\begin{proof}
تمرین.
\end{proof}

\begin{prob}
\olref[fol][syn][unq]{lem:no-prefix} را اثبات کنید.
\end{prob}

\begin{prop}
\ollabel{prop:unique-atomic}
اگر $!A$ یک !!{formula} اتمی باشد، یکی و فقط یکی از شرط‌های زیر را برآورده
می‌کند.
\begin{enumerate}
\tagitem{prvFalse}{$!A \ident \lfalse$.}{}
\tagitem{prvTrue}{$!A \ident \ltrue$.}{}
\item $!A \ident \Atom{R}{t_1,\dots,t_n}$، که در آن $R$ یک
  !!{predicate} $n$جایگاهی و $t_1$، \dots، $t_n$ ترم‌اند، و هر یک از
  $R$، $t_1$، \dots، $t_n$ به‌طور یکتا تعیین شده است.
\item $!A \ident \eq[t_1][t_2]$، که در آن $t_1$ و $t_2$ ترم‌هایی هستند
  که به‌طور یکتا تعیین شده‌اند.
\end{enumerate}
\end{prop}

\begin{proof}
تمرین.
\end{proof}

\begin{prob}
\olref[fol][syn][unq]{prop:unique-atomic} را اثبات کنید. (راهنمایی: نسخه‌ای
از \olref[fol][syn][unq]{lem:no-prefix} را برای ترم‌ها صورت‌بندی و اثبات
کنید.)
\end{prob}

\begin{prop}[خوانش یکتا]
هر !!{formula} یکی و فقط یکی از شرط‌های زیر را برآورده می‌کند.
\begin{enumerate}
\item $!A$ اتمی است.

\tagitem{prvNot}{$!A$ به صورت $\lnot !B$ است.}{}

\tagitem{prvAnd}{$!A$ به صورت $(!B \land !C)$ است.}{}

\tagitem{prvOr}{$!A$ به صورت $(!B \lor !C)$ است.}{}

\tagitem{prvIf}{$!A$ به صورت $(!B \lif !C)$ است.}{}

\tagitem{prvIff}{$!A$ به صورت $(!B \liff !C)$ است.}{}

\tagitem{prvAll}{$!A$ به صورت $\lforall[x][!B]$ است.}{}

\tagitem{prvEx}{$!A$ به صورت $\lexists[x][!B]$ است.}{}
\end{enumerate}
افزون بر این، در هر حالت $!B$، یا $!B$ و $!C$، به‌طور یکتا تعیین می‌شوند.
برای نمونه، این بدان معناست که دو جفت متفاوت $!B$، $!C$ و $!B'$، $!C'$
وجود ندارند که $!A$ هم به صورت
\iftag{prvIf}{$(!B \lif !C)$ و هم $(!B' \lif !C')$ باشد.}{
    \iftag{prvOr}{$(!B \lor !C)$ و هم $(!B' \lor !C')$ باشد.}{
      \iftag{prvAnd}{$(!B \land !C)$ و هم $(!B' \land !C')$ باشد.}{}}}
\end{prop}

\begin{proof}
قواعد ساخت ایجاب می‌کنند که اگر !!a{formula} اتمی نباشد، باید با یک
پرانتز باز~(، \iftag{prvNot}{$\lnot$،}{} یا یک سور آغاز شود. از سوی دیگر،
هر !!{formula} که با یکی از نمادهای زیر آغاز شود باید اتمی باشد:
!!a{predicate}، !!a{function}، !!a{constant}، !!a{variable}
\iftag{prvFalse}{، $\lfalse$}{}\iftag{prvTrue}{، $\ltrue$}{}.

پس در واقع تنها باید نشان دهیم که اگر $!A$ هم به صورت $(!B \ast
!C)$ و هم به صورت $(!B' \mathbin{\ast'} !C')$ باشد، آنگاه $!B \ident
!B'$، $!C \ident !C'$ و $\ast = {\ast'}$.

فرض کنید هم $!A \ident (!B \ast !C)$ و هم $!A \ident (!B'
\mathbin{\ast'} !C')$. آنگاه یا $!B \ident !B'$ است یا نیست. اگر باشد،
آشکار است که $\ast = {\ast'}$ و $!C \ident !C'$، زیرا در این صورت آن‌ها
زیررشته‌هایی از $!A$ هستند که از یک جایگاه آغاز می‌شوند و طول یکسانی دارند.
حالت دیگر $!B \not\ident !B'$ است. چون $!B$ و $!B'$ هر دو زیررشته‌هایی از
$!A$ هستند که از جایگاه واحدی آغاز می‌شوند، یکی باید پیشوند سرهٔ دیگری
باشد. اما بنا بر \olref{lem:no-prefix} این امر محال است.
\end{proof}


\end{document}
